\chapter{LITERATURE REVIEW}

This chapter reviews the role of sentiment in stock prediction, the language models and datasets used for financial sentiment analysis, and the research gaps addressed by this project.

\section{Stock Market Prediction and the Role of Sentiment}
The efficient market hypothesis holds that publicly available information, including news, is rapidly incorporated into asset prices \parencite{fama1970efficient}. This would limit the value of any single public signal, such as news sentiment. In practice, however, the incorporation of information may be neither instantaneous nor complete, and short-term reactions to news events may leave a window in which sentiment carries predictive value at short horizons. Media pessimism has been found to predict downward pressure on market prices, followed by a reversion towards fundamentals \parencite{tetlock2007giving}. At the firm level, the share of negative words in news about S\&P 500 companies predicted lower earnings, and prices briefly under-reacted to this information \parencite{tetlock2008more}. However, such effects have been found to be short-lived. In more than 900,000 news stories, daily news predicted returns for only about one to two days \parencite{heston2017news}. In about 12 million timestamped Dow Jones Newswires articles, the price response to news sentiment was largest on the day of publication, with a smaller response on the following day that had been fully absorbed within about three days \parencite{ke2019predicting}. This delay was longer and the response larger for small and volatile stocks, which is consistent with evidence that information spreads more slowly for smaller firms with less analyst coverage \parencite{hong2000bad}.
Industry surveys report sentiment among the most widely adopted categories of alternative data, but describe it as one input to be combined with prices, fundamentals or other data \parencite{aima2024casting}. The same framing was adopted in this project, in which sentiment was treated as a candidate feature to be tested for marginal predictive value and not as a standalone trading signal.

\section{FinBERT and Domain-Specific Language Models}
FinBERT is a BERT-based language model trained on financial text and fine-tuned for sentiment classification. It was developed because general sentiment models perform poorly on financial text, owing to domain-specific vocabulary and context-dependent polarity \parencite{araci2019finbert,yang2020finbert}. The most widely used benchmark is the Financial PhraseBank, a dataset of financial news sentences labelled as positive, negative or neutral by annotators with a financial market background, according to how a phrase might be expected to influence a stock price \parencite{malo2014good}. FinBERT and its variants have become the standard baseline against which newer methods, including prompted LLMs, are compared.

\section{Sentiment-Augmented Stock Prediction Models}
Recent studies have combined FinBERT sentiment with sequential and tree-based prediction models. A FinBERT-LSTM model was reported to improve on price-only LSTM and DNN baselines by weighting news according to whether it was market-wide, industry-specific or stock-specific \parencite{gu2024predicting}. Other work has applied FinBERT to social-media commentary, feeding the sentiment of Stocktwits posts into an ensemble Support Vector Machine for movement prediction \parencite{pmc2023stocktwits}. More recently, FinBERT sentiment has been combined with technical indicators in a gradient-boosted tree framework with SHAP-based explainability \parencite{mathematics2025stock}, and group-wise SHAP analysis has been used to assess the joint contribution of news keywords and technical indicators \parencite{arxiv2025groupshap}. These studies establish FinBERT-augmented prediction as an active research direction, but two gaps recur: evaluation on a single news source, and reported performance without a matched price-only baseline.

\section{General-Purpose LLMs versus Domain-Specific Models}
A parallel line of work has compared FinBERT with general-purpose LLMs on financial sentiment classification. As summarised in Table~\ref{tab:lit-llm}, the results are inconsistent and do not uniformly favour either approach.

\begin{longtable}{@{}P{4.3cm}P{3.6cm}P{2.3cm}P{3.0cm}@{}}
\caption{Reported FinBERT and General LLM Sentiment Accuracy}
\label{tab:lit-llm}\\
\toprule
\textbf{Study} & \textbf{Models compared} & \textbf{FinBERT} & \textbf{General LLM} \\
\midrule
\endfirsthead
\toprule
\textbf{Study} & \textbf{Models compared} & \textbf{FinBERT} & \textbf{General LLM} \\
\midrule
\endhead
\bottomrule
\endlastfoot
Forex news headlines \parencite{fatouros2023transforming} & FinBERT and prompted ChatGPT & 56.1\% & 73.0--79.0\%, across four prompts \\
NGX All-Share Index \parencite{ai2024innovative} & FinBERT, GPT-4 and logistic regression & 63.3\% & 54.2\% \\
Financial-sentiment benchmark of 5,843 sentences \parencite{vidhya2025finetuning} & FinBERT and fine-tuned GPT-4o mini & 75.8\% & 76.5\% \\
Five news sources \parencite{liu2025enhancing} & FinBERT and GPT-2 & 32.2--75.6\%, by source & 34.4--65.5\%, by source \\
New York Times news by sector \parencite{electronics2025comparative} & FinBERT and prompted GPT-4o & Lower than GPT-4o & Up to 10\% higher, by sector \\
Market reaction to Bloomberg headlines \parencite{arxiv2024optimizing} & FinBERT, DistilRoBERTa and zero-shot GPT-4 & Weighted F1 0.44 & Weighted F1 0.47 \\
\end{longtable}

Three observations follow from this comparison. First, general LLMs have outperformed FinBERT when carefully prompted or fine-tuned, but not consistently when used with a fixed prompt. Second, the accuracy of FinBERT has been highly dependent on the source, varying by more than 40 percentage points across news outlets in a single study, which motivated the multi-source evaluation in this project. Third, several of the studies favouring general LLMs relied on proprietary, API-based models, which are costly at the scale of a multi-year backtest and are not fully reproducible, as the underlying model may change.

For these reasons, and following the guidance of the supervisor, FinBERT was not benchmarked against proprietary or agentic AI systems. A comparison was instead made with open, versioned, non-agentic models, including a small LLM used with one fixed prompt, so that every checkpoint could be reproduced.

\section{Financial Sentiment and Stock Price Datasets}
The Financial PhraseBank remains the standard labelled dataset for financial sentiment classification \parencite{malo2014good}. For paired news and price data, the Financial News and Stock Price Integration Dataset (FNSPID) is a large public resource, combining approximately 29.7 million daily stock price records with 15.7 million time-aligned news articles for 4,775 S\&P~500 companies between 1999 and 2023. Its authors reported that sentiment information provided a modest additional improvement in model performance \parencite{dong2024fnspid}.

\section{Research Gap}
The reviewed literature shows that FinBERT-augmented stock prediction is technically viable and well studied, and that the choice between a domain-specific model and a general-purpose LLM remains unsettled. What existing studies do not consistently provide is a controlled comparison of FinBERT across several news sources within one study, and a rigorous ablation that isolates the contribution of sentiment over a matched price-only baseline across several horizons. As shown in Table~\ref{tab:lit-gap}, this project was designed to address both gaps.

\begin{utptable}{Comparison of Related Work with This Project}
\label{tab:lit-gap}
\begin{tabular}{@{}P{6.2cm}P{3.2cm}P{1.8cm}P{1.8cm}@{}}
\toprule
\textbf{Related work} & \textbf{Sentiment model} & \textbf{Multiple sources} & \textbf{Price-only ablation} \\
\midrule
FinBERT-LSTM \parencite{gu2024predicting} & FinBERT & No & Partial \\
FinBERT, GPT-4 and logistic regression \parencite{ai2024innovative} & FinBERT and GPT-4 & No & No \\
FinBERT, XGBoost and SHAP \parencite{mathematics2025stock} & FinBERT & No & No \\
This project & FinBERT, with six comparison models & Yes & Yes \\
\bottomrule
\end{tabular}
\end{utptable}
